% Candidate Figure 2. Four-panel, transparent, formula-free mechanism view.
\begin{figure}[H]
\centering
\begin{tikzpicture}[
  every node/.style={font=\fontsize{8.2}{9.6}\selectfont},
  panel/.style={rectangle,draw=black,line width=0.45pt,align=center,
    inner sep=2.6mm},
  plan/.style={rectangle,draw=black,line width=0.38pt,minimum width=17mm,
    minimum height=5.5mm,align=center,inner sep=1pt},
  route/.style={rectangle,draw=black,line width=0.38pt,minimum width=43mm,
    minimum height=5.7mm,align=center,inner sep=1pt},
  selected/.style={route,line width=0.85pt},
  rejected/.style={plan,dashed},
  action/.style={rectangle,draw=black,line width=0.45pt,text width=45mm,
    minimum height=8.2mm,align=center,inner sep=1.5pt},
  arrow/.style={-{Latex[length=2.5mm,width=1.65mm]},line width=0.55pt,
    line cap=rect,line join=miter,shorten >=1.2pt},
  small/.style={font=\fontsize{7.9}{9.1}\selectfont,align=center,inner sep=0.8pt}
]
% Panel (a): source archives.
\node[panel,minimum width=66mm,minimum height=44mm] (pa) at (-3.60,2.70) {};
\node[small] at (-3.60,4.52) {(a) 三视角候选档案};
\node[small] at (-5.55,3.86) {HGS-F};
\node[small] at (-3.60,3.86) {HGS-E};
\node[small] at (-1.65,3.86) {HGS-M};
\foreach \x in {-5.55,-3.60,-1.65} {
  \node[plan] at (\x,3.28) {候选方案1};
  \node[plan] at (\x,2.58) {候选方案2};
  \node[plan] at (\x,1.88) {候选方案3};
  \node[small] at (\x,1.30) {$\vdots$};
}

% Panel (b): full-model review.
\node[panel,minimum width=51mm,minimum height=44mm] (pb) at (3.85,2.70) {};
\node[small] at (3.85,4.52) {(b) 完整模型统一检查};
\node[action] at (3.85,3.55)
  {车型指派、补能、时间窗\\车队数量、分时成本与碳排放};
\node[plan,minimum width=20mm] at (2.55,2.36) {可行候选};
\node[plan,minimum width=20mm] at (5.15,2.36) {可行候选};
\node[rejected,minimum width=20mm] at (2.55,1.52) {不可行剔除};
\node[rejected,minimum width=20mm] at (5.15,1.52) {不可行剔除};

% Panel (c): route pool.
\node[panel,minimum width=51mm,minimum height=44mm] (pc) at (-4.35,-2.70) {};
\node[small] at (-4.35,-0.88) {(c) 可行路线合并去重};
\node[route] at (-4.35,-1.62) {广州场--客户4--客户18--广州场};
\node[selected] at (-4.35,-2.39) {深圳场--客户7--客户21--深圳场};
\node[route] at (-4.35,-3.16) {广州场--客户2--客户13--广州场};
\node[selected] at (-4.35,-3.93) {深圳场--客户5--客户16--深圳场};

% Panel (d): MIP and final adjudication.
\node[panel,minimum width=66mm,minimum height=44mm] (pd) at (3.10,-2.70) {};
\node[small] at (3.10,-0.88) {(d) 限时组合与最终择优};
\node[action,text width=54mm] (mip) at (3.10,-1.62)
  {限时MIP从路线池选取\\客户恰好覆盖的可行组合};
\node[action,text width=54mm] (check) at (3.10,-2.78)
  {组合方案重新完成排班\\并通过完整模型与独立检查};
\node[action,text width=54mm] (best) at (3.10,-3.94)
  {与三视角最好完整解比较\\输出成本最低的可行方案};

% Orthogonal arrows, routed only through panel margins.
\draw[arrow] (pa.east) -- (pb.west);
\draw[arrow] (pb.south) -- (3.85,0.08) -- (-4.35,0.08) -- (pc.north);
\draw[arrow] (pc.east) -- (pd.west);
\draw[arrow] (mip) -- (check);
\draw[arrow] (check) -- (best);
\draw[arrow] (pb.east) -- (6.95,2.70) -- (6.95,-3.94) -- (best.east);
\node[small,above=0.7mm] at (0.18,2.70) {完整检查};
\node[small,above=0.7mm] at (-0.62,0.08) {可行路线进入池};
\node[small,above=0.7mm] at (-0.58,-2.70) {路线池};
\node[small,text width=17mm,anchor=east] at (6.78,-0.24)
  {最好完整解直接比较};
\end{tikzpicture}
\caption{多视角候选的完整模型筛选与路线重组}
\label{fig:route-pool}
\end{figure}
